\documentclass[twocolumn]{article}
\usepackage[T1]{fontenc}
\usepackage{lmodern,amsmath,amsthm,booktabs}
\newtheorem{theorem}{Theorem}
\begin{document}

\section{Background}

Earlier work by Smith~\cite{smith2020} and by Jones~\cite[p.~4]{jones2019} measured the effect, and the effect was confirmed again in Section~\ref{sec:method}, where the effect is defined. Schrödinger's café was naïve about the effect, and the effect of the café on Schrödinger was small.

A \emph{very} important result is that results repeat: the result of one run predicts the result of the next run, and every run has a result. Each run is short, and the next run is shorter than the run before it.

\section{Method}\label{sec:method}

\begin{theorem}[Main result]
Every bounded sequence has a convergent subsequence, and every convergent sequence is bounded.
\end{theorem}

\begin{proof}
Take a bounded sequence and split the interval in half; one half holds infinitely many terms of the sequence.
\end{proof}

\begin{table}[t]
  \centering
  \begin{tabular}{@{}lrr@{}}
    \toprule
    \multicolumn{3}{c}{Measured values} \\
    \midrule
    Sample & Mass & Width \\
    alpha & 12.5 & 3.1 \\
    beta & 13.0 & 2.9 \\
    \bottomrule
  \end{tabular}
  \caption{Values for every sample, with the mass and the width of each sample.}
  \label{tab:values}
\end{table}

\begin{description}
  \item[Mass] The mass of the sample, in grams.
  \item[Width] The width of the sample, in millimetres.
\end{description}

\begin{quote}
A quoted passage about the sample, set apart from the text around the sample.
\end{quote}

Table~\ref{tab:values} lists the values; the mass of each sample is in the second column and the width of each sample is in the third.

\begin{thebibliography}{9}
\bibitem{smith2020} A. Smith, \emph{A study of effects}, 2020.
\bibitem{jones2019} B. Jones, \emph{More effects}, 2019.
\end{thebibliography}

\end{document}
